\documentclass[11pt]{article}

\usepackage[T1]{fontenc}
\usepackage{lmodern}
\usepackage[margin=1in]{geometry}
\usepackage{amsmath,amssymb,amsthm,mathtools}
\usepackage{booktabs}
\usepackage{array}
\usepackage{enumitem}
\usepackage{hyperref}
\usepackage{microtype}

\hypersetup{
  colorlinks=true,
  linkcolor=blue,
  urlcolor=blue,
  pdftitle={Geodesic Kamada--Kawai implementation plan for grip},
  pdfauthor={Codex}
}

\setlength{\parskip}{0.6em}
\setlength{\parindent}{0pt}

\newcommand{\RR}{\mathbb{R}}
\newcommand{\norm}[1]{\left\lVert #1 \right\rVert}
\newcommand{\vect}[1]{\mathbf{#1}}

\title{Geodesic Kamada--Kawai in \texttt{grip}\\
\large Algorithm Details and Implementation Plan}
\author{Generated from direct inspection of the current \texttt{grip} source tree}
\date{\today}

\begin{document}

\maketitle

\begin{abstract}
This note turns the geodesic Kamada--Kawai (GKK) functional into a concrete
implementation plan for \texttt{grip}. The package already has landmark geodesic
Kamada--Kawai (LGKK) scoring and optimization in both R and C++. The missing
piece is the full all-pairs GKK variant so that we can compare three objects on
the same graphs: a classical KK objective based on chord length, a full GKK
objective based on path length along the embedded graph, and the sparse LGKK
approximation already present in the package. The main conclusion is that GKK
fits naturally into the current codebase: once the all-pairs path cache is
available, almost all of the existing LGKK path-length, scale-fitting, energy,
and gradient machinery can be reused with only light refactoring.
\end{abstract}

\section{Goal}

The target deliverable is a full all-pairs geodesic Kamada--Kawai implementation
for connected weighted or unweighted graphs, together with a practical plan for
an R reference implementation and a faster C++ backend with an R interface.

The immediate purpose is comparison. We want to quantify when replacing the
classical KK chord-length discrepancy by an embedded-path discrepancy actually
improves preservation of graph geodesics, and what runtime cost is paid for that
gain.

The document therefore focuses on four questions:
\begin{enumerate}[leftmargin=1.5em]
\item What exactly is the GKK objective that \texttt{grip} should implement?
\item How does that objective relate to the current LGKK code?
\item What is the simplest implementation path that gives us a correct baseline
first and a fast backend second?
\item Which design decisions should be made explicitly before coding?
\end{enumerate}

\section{Algorithmic Definition}

\subsection{Graph, embedding, and chosen shortest paths}

Let $G = (V,E)$ be a connected graph with vertex set
$V = \{1,\dots,n\}$. Let the graph distance between vertices $i$ and $j$ be
$g_{ij}$. In the unweighted case this is hop count; in the weighted case it is
the shortest-path distance induced by the positive edge weights.

For each unordered pair $i < j$, choose one deterministic graph shortest path
\[
\gamma_{ij} = (v_0 = i, v_1, \dots, v_r = j).
\]
The current LGKK code already uses this model: it does not average over many
shortest paths and it does not recompute the chosen path from the current
embedding. The most natural full GKK implementation should keep that same rule.

Given coordinates
$X = (\vect{x}_1,\dots,\vect{x}_n)$ with $\vect{x}_i \in \RR^d$,
define the embedded length of an edge $e = \{u,v\}$ by
\[
\lambda_e(X) = \sqrt{\norm{\vect{x}_u - \vect{x}_v}^2 + \varepsilon_{\mathrm{edge}}^2},
\]
where $\varepsilon_{\mathrm{edge}} \ge 0$ is the same small stabilizer already
used by the LGKK code.

The embedded geodesic length of the chosen shortest path $\gamma_{ij}$ is
\[
h_{ij}(X) = \sum_{e \in \gamma_{ij}} \lambda_e(X).
\]

\subsection{Full GKK energy}

With global stiffness constant $K > 0$, graph-distance floor
$\varepsilon_{\mathrm{dist}} > 0$, and scale factor $L_0 > 0$, define
\[
k_{ij} = \frac{K}{\max(g_{ij}, \varepsilon_{\mathrm{dist}})^2}
\]
and
\[
E_{\mathrm{GKK}}(X;L_0)
=
\frac{1}{2}
\sum_{1 \le i < j \le n}
k_{ij}\bigl(h_{ij}(X) - L_0 g_{ij}\bigr)^2.
\]

This is the direct geodesic analogue of the classical KK energy. The only
structural change is that the discrepancy term uses path length measured along
the current embedded graph rather than the direct Euclidean chord
$\norm{\vect{x}_i-\vect{x}_j}$.

\subsection{Analytic scale fitting}

For scoring and for optimization with a scale-invariant objective, it is useful
to fit $L_0$ analytically for the current layout. For fixed $X$,
\[
L_0^\star(X)
=
\operatorname*{arg\,min}_{L_0} E_{\mathrm{GKK}}(X;L_0)
=
\frac{\sum_{i<j} k_{ij} g_{ij} h_{ij}(X)}
{\sum_{i<j} k_{ij} g_{ij}^2}.
\]

This is exactly the same algebra used in the existing LGKK scorer.

The profiled objective is
\[
\widetilde{E}_{\mathrm{GKK}}(X) = E_{\mathrm{GKK}}(X;L_0^\star(X)).
\]

An important implementation fact is that the gradient of the profiled objective
is still easy to compute: at $L_0^\star(X)$ we have
$\partial E_{\mathrm{GKK}} / \partial L_0 = 0$, so the total derivative of
$\widetilde{E}_{\mathrm{GKK}}$ is obtained by differentiating
$E_{\mathrm{GKK}}(X;L_0)$ with $L_0$ held fixed at the optimal value. In
practice this means we can recompute $L_0^\star$ at every evaluation without
adding a derivative term for $dL_0^\star / dX$.

\subsection{Gradient}

For one pair $(i,j)$ with residual
\[
r_{ij}(X;L_0) = h_{ij}(X) - L_0 g_{ij},
\]
the pair contribution to the energy is
\[
\frac{1}{2}k_{ij} r_{ij}(X;L_0)^2.
\]

If an oriented edge $(u,v)$ belongs to $\gamma_{ij}$, then
\[
\nabla_{\vect{x}_u}\lambda_{\{u,v\}}(X)
=
\frac{\vect{x}_u - \vect{x}_v}{\lambda_{\{u,v\}}(X)},
\qquad
\nabla_{\vect{x}_v}\lambda_{\{u,v\}}(X)
=
-
\frac{\vect{x}_u - \vect{x}_v}{\lambda_{\{u,v\}}(X)}.
\]

Hence
\[
\nabla_{\vect{x}_u} E_{\mathrm{GKK}}(X;L_0)
=
\sum_{i<j}
k_{ij} r_{ij}(X;L_0)
\sum_{e \in \gamma_{ij} \,:\, u \in e}
\sigma(u,e)
\frac{\vect{x}_{a(e)} - \vect{x}_{b(e)}}{\lambda_e(X)},
\]
where $\sigma(u,e) = +1$ if $u = a(e)$ and $\sigma(u,e) = -1$ if $u = b(e)$.

This is the same pattern already implemented in \texttt{grip.lgkk.energy.gradient()}:
\begin{enumerate}[leftmargin=1.5em]
\item compute the embedded length of each edge on the chosen path,
\item sum those edge lengths to get $h_{ij}(X)$,
\item form the residual,
\item distribute the residual back along the path edges.
\end{enumerate}

\subsection{Complexity}

Define the total number of stored path-edge references by
\[
M = \sum_{1 \le i < j \le n} \bigl(|\gamma_{ij}| - 1\bigr).
\]

Then:
\begin{itemize}[leftmargin=1.5em]
\item all-pairs graph-distance computation costs
$O(nm)$ for unweighted graphs with BFS from every source, or
$O(nm \log n)$ for weighted graphs with Dijkstra from every source;
\item storing the distance matrix costs $O(n^2)$ memory;
\item storing all chosen paths explicitly costs $O(M)$ memory;
\item one full energy or gradient evaluation costs $O(Md)$ arithmetic.
\end{itemize}

This is the main practical difference between GKK and LGKK. For the sparse LGKK
variant, the pair count is closer to
$O(n(\texttt{local\_nbrs} + \texttt{landmark\_count}))$; for full GKK it is
$n(n-1)/2$, and on long-diameter graphs $M$ can be cubic in $n$.

\section{What \texttt{grip} Already Has}

Inspection of the current source tree shows that \texttt{grip} already contains
most of the mathematical infrastructure needed by GKK.

\subsection{R-level LGKK pieces}

The key functions in \texttt{R/grip\_quality.R} are:
\begin{itemize}[leftmargin=1.5em]
\item \texttt{grip.prepare.landmark.geodesic.kk()}
\item \texttt{grip.score.landmark.geodesic.kk()}
\item \texttt{grip.optimize.landmark.geodesic.kk()}
\item \texttt{grip.lgkk.path.lengths()}
\item \texttt{grip.lgkk.fit.scale()}
\item \texttt{grip.lgkk.energy.gradient()}
\end{itemize}

These routines already do the following:
\begin{enumerate}[leftmargin=1.5em]
\item validate weighted and unweighted undirected graphs;
\item sort adjacency for deterministic tie-breaking;
\item compute deterministic shortest-path trees from every source;
\item reconstruct a chosen shortest path for each retained pair;
\item measure path length along the embedded graph;
\item fit the global scale analytically;
\item evaluate the geodesic energy and its gradient;
\item optimize by deterministic gradient descent with backtracking.
\end{enumerate}

\subsection{Compiled LGKK pieces}

The compiled multiscale LGKK code lives primarily in:
\begin{itemize}[leftmargin=1.5em]
\item \texttt{src/DrawGraph.h}
\item \texttt{src/MishSupport.cpp}
\item \texttt{src/grip\_rcpp.cpp}
\end{itemize}

At each active MISF level, the current implementation:
\begin{enumerate}[leftmargin=1.5em]
\item computes active-set shortest-path distances,
\item selects sparse local and landmark targets,
\item reconstructs chosen shortest paths,
\item fits an initial scale,
\item runs a gradient-descent refinement over the sparse pair set.
\end{enumerate}

\subsection{Most important observation}

The current LGKK scorer and optimizer are already generic in everything that
matters except pair selection. Once a prepared object contains:
\begin{itemize}[leftmargin=1.5em]
\item a pair matrix,
\item graph distances for those pairs,
\item path vertices or path edges for those pairs,
\end{itemize}
the downstream code does not really care whether the pair set came from
landmark sparsification or from the full all-pairs list.

That means full GKK should be implemented by changing the preparation layer
first, then lightly refactoring the shared energy machinery so LGKK and GKK
call the same internal core.

\section{Recommended Public API}

The cleanest user-facing API is the full analogue of the existing LGKK API:
\begin{align*}
&\texttt{grip.prepare.geodesic.kk(...)} \\
&\texttt{grip.score.geodesic.kk(...)} \\
&\texttt{grip.optimize.geodesic.kk(...)}
\end{align*}

The prepared object should have class \texttt{"grip\_gkk\_prepared"} and should
reuse the same broad schema as \texttt{"grip\_lgkk\_prepared"}:
\begin{itemize}[leftmargin=1.5em]
\item \texttt{n}
\item \texttt{edges}
\item \texttt{adj\_list}
\item \texttt{weight\_list}
\item \texttt{pair\_matrix}
\item \texttt{pair\_graph\_distance}
\item \texttt{path\_vertices}
\item \texttt{path\_edges}
\item \texttt{graph\_diameter}
\item \texttt{distance\_matrix}
\end{itemize}

The difference from LGKK is that \texttt{pair\_matrix} is now the full all-pairs
set
\[
\{(i,j) : 1 \le i < j \le n\}.
\]

For scoring and optimization, the arguments should stay as close as possible to
the existing LGKK names so that the comparison code is easy to write and easy to
read. In particular, the following arguments can be retained with identical
meaning:
\begin{itemize}[leftmargin=1.5em]
\item \texttt{stiffness}
\item \texttt{distance\_floor}
\item \texttt{edge\_length\_epsilon}
\item \texttt{max\_iter}
\item \texttt{initial\_step}
\item \texttt{step\_shrink}
\item \texttt{armijo\_factor}
\item \texttt{grad\_tol}
\item \texttt{min\_step}
\item \texttt{recenter}
\item \texttt{return\_trace}
\end{itemize}

I also recommend adding one new scale-policy argument, for example
\texttt{scale\_mode = c("profiled", "fixed\_initial", "user")}, because scale
handling is one of the main choices that affects comparability.

\section{Implementation Strategy}

\subsection{Core idea}

The implementation should be organized around a shared geodesic KK core and two
pair-selection front ends:
\begin{itemize}[leftmargin=1.5em]
\item LGKK: sparse local-plus-landmark pair selection,
\item GKK: full all-pairs selection.
\end{itemize}

This avoids copying the energy and gradient logic into a second code path.

\subsection{Phase 1: R reference implementation}

The fastest route to a correct baseline is an R implementation in
\texttt{R/grip\_quality.R}. The main steps are:

\begin{enumerate}[leftmargin=1.5em]
\item Refactor the current LGKK preparation code into a reusable internal helper
that computes sorted adjacency, all-source shortest-path trees, the full
distance matrix, and parent arrays.

\item Refactor the current downstream helpers so they are pair-set agnostic.
Concretely, the current LGKK helpers for path lengths, analytic scale fitting,
and energy-gradient evaluation should either be renamed to generic internal
helpers or wrapped around a generic geodesic KK core.

\item Implement \texttt{grip.prepare.geodesic.kk()} by constructing the all-pairs
\texttt{pair\_matrix} and reconstructing one chosen path for every pair from the
precomputed parent arrays.

\item Implement \texttt{grip.score.geodesic.kk()} on top of the shared scorer.
This should report the same style of summary as the LGKK scorer so that
benchmark tables line up naturally.

\item Implement \texttt{grip.optimize.geodesic.kk()} by reusing the current
deterministic gradient-descent loop. The first optimizer does not need a new
algorithm; it only needs the full pair set.
\end{enumerate}

This phase should give us a correct, readable, testable GKK reference
implementation quickly.

\subsection{Phase 2: compiled backend}

The compiled backend should accelerate repeated energy and gradient evaluation,
which is the true bottleneck during optimization. I do \emph{not} recommend
putting the full all-pairs GKK implementation directly into the current
\texttt{DrawGraph} multiscale engine as the first compiled step. That engine is
organized around active-set refinement, while full GKK is a standalone
all-pairs objective.

Instead, the cleaner route is:
\begin{enumerate}[leftmargin=1.5em]
\item add a standalone C++ module, for example
\texttt{src/GeodesicKk.h} and \texttt{src/GeodesicKk.cpp};
\item expose Rcpp entry points from \texttt{src/grip\_rcpp.cpp};
\item keep the public R API in R, with the C++ path hidden behind the same
wrappers.
\end{enumerate}

For C++ efficiency, the prepared path cache should be flattened. Rather than a
list of edge matrices, use arrays such as:
\begin{itemize}[leftmargin=1.5em]
\item \texttt{pair\_offset}
\item \texttt{path\_edge\_u}
\item \texttt{path\_edge\_v}
\item \texttt{pair\_graph\_distance}
\end{itemize}
so that the compiled evaluator can stream through contiguous memory.

The C++ backend can be introduced in two layers:
\begin{enumerate}[leftmargin=1.5em]
\item compiled scoring and gradient from an R-prepared flattened cache;
\item later, compiled preparation if the all-source shortest-path stage becomes
the dominant cost on the graphs we care about.
\end{enumerate}

\subsection{Phase 3: comparison tooling}

Once the R reference and compiled backend exist, add benchmark and comparison
tools that evaluate at least:
\begin{itemize}[leftmargin=1.5em]
\item classical KK or the nearest available full KK reference,
\item full GKK,
\item sparse LGKK.
\end{itemize}

The comparison scripts should use the same starting coordinates, the same
dimension, and the same stopping policy whenever possible. Otherwise differences
in initialization may dominate differences in objective.

\section{Concrete File-Level Plan}

\begin{center}
\begin{tabular}{>{\raggedright\arraybackslash}p{0.28\linewidth} >{\raggedright\arraybackslash}p{0.62\linewidth}}
\toprule
Location & Planned work \\
\midrule
\texttt{R/grip\_quality.R} &
Add \texttt{grip.prepare.geodesic.kk()},
\texttt{grip.score.geodesic.kk()},
\texttt{grip.optimize.geodesic.kk()}, and refactor the current LGKK helpers into
generic geodesic KK internals shared by LGKK and GKK. \\
\addlinespace
\texttt{man/} &
Generate Rd files for the three new public functions. \\
\addlinespace
\texttt{tests/testthat/}\newline\texttt{test-geodesic-kk.R} &
Add exactness, determinism, optimizer-decrease, weighted-graph, and R-vs-C++
agreement tests. \\
\addlinespace
\texttt{src/GeodesicKk.h}, \texttt{src/GeodesicKk.cpp} &
Add standalone compiled scorer and optimizer helpers for the flattened all-pairs
path cache. \\
\addlinespace
\texttt{src/grip\_rcpp.cpp}, \texttt{src/RcppExports.cpp}, \texttt{R/RcppExports.R} &
Expose compiled helpers through Rcpp and wire them into the R wrappers. \\
\addlinespace
\texttt{tools/benchmarks/gkk\_lgkk\_paper/*.R} &
Add timing and quality benchmarks comparing KK, GKK, and LGKK on small and
medium graphs. \\
\bottomrule
\end{tabular}
\end{center}

\section{Design Decisions to Make Explicitly}

\subsection*{1. Path-choice rule}

\textbf{Question:} for each pair $(i,j)$, do we use one canonical shortest path
or somehow combine multiple shortest paths?

\textbf{Recommendation:} use one deterministic canonical shortest path, exactly
as LGKK already does.

\textbf{Reason:} it keeps the functional deterministic, keeps the prepared cache
compact, matches the current code style, and avoids a much heavier objective.

\subsection*{2. Scale handling}

\textbf{Question:} should $L_0$ be profiled at every evaluation, fitted once at
the initial layout, or supplied by the user?

\textbf{Recommendation:} support all three modes, with \texttt{"profiled"} as the
default for scoring and probably also for the final optimizer.

\textbf{Reason:} profiled scale is mathematically clean, directly comparable
across layouts, and does not complicate the gradient. The current LGKK optimizer
fixes the initial scale for robustness, but GKK is a good opportunity to expose
the scale policy explicitly.

\subsection*{3. Optimizer family}

\textbf{Question:} should GKK use a classical KK-style per-vertex Newton method,
the current whole-layout gradient descent, or a quasi-Newton method such as
L-BFGS?

\textbf{Recommendation:} start with the current deterministic gradient descent.

\textbf{Reason:} the existing LGKK optimizer already matches the geodesic path
objective, is straightforward to validate, and is enough for a first correct
implementation. A later L-BFGS upgrade can be added if convergence is too slow.

\subsection*{4. Shared-core refactor versus duplicate code}

\textbf{Question:} do we add separate GKK-specific energy code, or do we
genericize the existing LGKK internals?

\textbf{Recommendation:} genericize the existing internals.

\textbf{Reason:} the mathematical core is already pair-set agnostic. Duplicating
it would create unnecessary maintenance risk and make R-vs-LGKK comparisons less
trustworthy.

\subsection*{5. Prepared-object representation}

\textbf{Question:} should the prepared object store only R-friendly lists of
paths, or should it also store a flattened C++-friendly cache?

\textbf{Recommendation:} keep the user-visible list structure, but allow optional
flattened arrays inside the prepared object for compiled evaluation.

\textbf{Reason:} this preserves readability at the R level and gives the C++
backend contiguous storage without introducing external-pointer complexity on day
one.

\subsection*{6. Size guardrails}

\textbf{Question:} what should happen on graphs where all-pairs GKK is too large?

\textbf{Recommendation:} warn or error explicitly once pair count or total
path-edge references exceed a configurable threshold; do not silently fall back
to LGKK.

\textbf{Reason:} silent fallback would make the comparison experiments ambiguous.

\subsection*{7. Weighted and disconnected graphs}

\textbf{Question:} which graph classes should the first version support?

\textbf{Recommendation:} support weighted and unweighted connected graphs from
the beginning; reject disconnected graphs initially, mirroring the current LGKK
behavior.

\textbf{Reason:} weighted support is already present in the current helper stack,
while disconnected-graph semantics would require extra policy decisions about
infinite distances and component placement.

\subsection*{8. Standalone tool versus multiscale integration}

\textbf{Question:} should full GKK be a standalone scorer/optimizer first, or a
new multiscale refinement stage inside \texttt{grip.layout.globalrep()}?

\textbf{Recommendation:} standalone first.

\textbf{Reason:} the user need is comparison with KK and LGKK, not immediate
integration into the production GRIP layout pipeline. A standalone tool is much
easier to validate and benchmark.

\subsection*{9. Comparison target for ``pure KK''}

\textbf{Question:} what exactly is the KK baseline?

\textbf{Recommendation:} decide this explicitly before benchmarking. Either
implement a matching full classical KK reference in \texttt{grip} or use one
external reference implementation with controlled initialization.

\textbf{Reason:} the current package core is centered on GRIP-style local and
multiscale routines rather than a textbook all-pairs KK solver, so the baseline
must be pinned down carefully to make the GKK comparison meaningful.

\section{Testing Plan}

The first test file should mirror the structure of the current LGKK tests.
Recommended tests are:
\begin{enumerate}[leftmargin=1.5em]
\item \textbf{Exact weighted path realization.} A path graph embedded with edge
lengths proportional to the graph weights should give zero or near-zero GKK
energy after scale fitting.

\item \textbf{Deterministic preparation.} Repeated calls to
\texttt{grip.prepare.geodesic.kk()} on the same graph should return identical
pair matrices, distances, and chosen paths.

\item \textbf{Scoring detects perturbation.} A canonical geometric graph such as
a grid or Sierpinski-carpet graph should score better than a noisy perturbation
under GKK.

\item \textbf{Optimizer decreases objective.} The optimizer should reduce GKK
energy and relative path error from a perturbed start.

\item \textbf{R and C++ agreement.} The compiled backend should match the R
reference to tight numeric tolerance on both energy and gradient-derived updates.

\item \textbf{Weighted graph support.} Weighted graphs should be covered in both
preparation and optimization tests.

\item \textbf{Guardrail behavior.} Oversized graphs should emit the documented
warning or error rather than failing later in the loop.
\end{enumerate}

\section{Benchmark Plan}

To evaluate the advantage of GKK over classical KK and the approximation quality
of LGKK, benchmark at least the following graph families:
\begin{itemize}[leftmargin=1.5em]
\item paths and ladders,
\item rectangular grids,
\item trees with long diameters,
\item Sierpinski-carpet or other fractal-like grids,
\item weighted road-like or corridor-like graphs,
\item cycles with shortcuts or chord edges.
\end{itemize}

For each method, record:
\begin{itemize}[leftmargin=1.5em]
\item runtime,
\item iteration count,
\item final classical KK energy if available,
\item final GKK energy,
\item geodesic weighted RMSE and relative RMSE,
\item sampled stress or another chord-based distortion metric,
\item visual examples for a small representative subset.
\end{itemize}

This benchmark matrix is important because GKK is not expected to dominate KK on
every metric. The hypothesis is narrower: GKK should better preserve graph
geodesic structure, especially on graphs where a straight chord is a poor proxy
for travel along the graph.

\section{Recommended Order of Work}

\begin{enumerate}[leftmargin=1.5em]
\item Build the R reference implementation by reusing the current LGKK prepared
schema and genericizing the scorer/gradient helpers.

\item Add tests until the R version is trustworthy on exact, weighted, and
perturbed examples.

\item Add benchmark scripts and obtain first comparison results on small graphs.

\item Implement the compiled scorer and optimizer over a flattened prepared cache.

\item Compare R and C++ numerically, then rerun the benchmarks.

\item Only after that, decide whether GKK should be integrated into the main
global-rep layout pipeline or kept as a standalone comparison and polishing
tool.
\end{enumerate}

\section{Conclusion}

The full geodesic Kamada--Kawai method is a natural next step for
\texttt{grip}. The package already contains almost everything needed:
deterministic shortest-path preparation, path-based energy evaluation, analytic
scale fitting, and a working geodesic optimizer. The missing part is mainly the
all-pairs preparation layer and a small refactor that turns the existing LGKK
machinery into a shared geodesic KK core.

The most important practical message is this: GKK should be implemented first as
a standalone all-pairs scorer and optimizer, not as an immediate extension of
the multiscale engine. That route gives the cleanest comparison against KK and
LGKK, keeps the code understandable, and lets us decide later whether the extra
cost of full GKK is justified by measurable gains in geodesic preservation.

\end{document}
